% ============================================================
%  GUÍA — DERIVADAS APLICADAS A LA ECONOMÍA
%  Matemática Aplicada a la Economía · Temas 1–4
%  Fuente: tex_files/clase_derivadas_economía_genesis/clases.md
%  Autor: Ing. Gabriel Astudillo
%  Compilar:  latexmk -xelatex guia_derivadas_economia.tex
%  Nota: la definición de derivada se escribe con \Delta x.
% ============================================================

\documentclass[12pt, letterpaper]{article}

% ── Paquetes ─────────────────────────────────────────────────
\usepackage{mathwizards-guia}
\mwguia{Guía — Derivadas Aplicadas a la Economía}{Ing. Gabriel Astudillo $\cdot$ Matemática Aplicada a la Economía}

\begin{document}
% ============================================================

% ── Portada / Título ─────────────────────────────────────────
\mwportada{Derivadas Aplicadas a la Economía}{Variación, modelos económicos y cálculo diferencial}{De la tasa de cambio a la derivada}

\vspace{4pt}
\begin{cuadroinfo}
  \small
  \textbf{Presentación:} Esta guía reúne el contenido del curso \emph{Matemática Aplicada a la Economía}. Partimos del \textbf{incremento} de una variable y de la \textbf{tasa de cambio} (Tema 1); construimos los \textbf{modelos económicos} clásicos —ingreso, costo, beneficio, equilibrio de mercado e interés compuesto— (Tema 2); formalizamos la \textbf{derivada} como límite y su interpretación geométrica (Tema 3); y cerramos con las \textbf{reglas y técnicas de derivación}: producto, cociente, cadena, exponencial, logaritmo, derivación logarítmica, implícita y sucesivas (Tema 4).
  \\[6pt]
  Los \textbf{Ejemplos Resueltos} provienen de las clases; los \textbf{Ejercicios Propuestos} incluyen versiones nuevas y los ejercicios de clase aún sin desarrollar, marcados con \textbf{(Clase)}.
  \\[6pt]
  \textbf{Notación clave:} $\Delta x$ (incremento de $x$), $\Delta y$ (incremento de $y$), $T=\Delta y/\Delta x$ (tasa de cambio media).
  \\[6pt]
  \textbf{Nivel de Dificultad:}\quad
  \medio{} Intermedio $(\bigstar\bigstar)$ \quad
  \dificil{} Difícil $(\bigstar\bigstar\bigstar)$ \quad
  \muyDificil{} Muy difícil $(\bigstar\bigstar\bigstar\bigstar)$
\end{cuadroinfo}

\vspace{8pt}

% ============================================================
\section{Marco Teórico}
% ============================================================

% ─────────────────────────────────────────────────────────────
\subsection{Tema 1 · Variación e incremento de una función}

\begin{cuadroteoria}
  \textbf{Incremento de la variable independiente.} Se denota $\Delta x$ (``delta'' $=$ incremento) y se define como la diferencia entre dos valores de $x$:
  \[
    \Delta x = x_2 - x_1 = x_f - x_0
    \qquad\Longrightarrow\qquad
    x_2 = x_1 + \Delta x
  \]
  \textbf{Incremento de la función.} Se denota $\Delta y$ o $\Delta F$:
  \[
    \Delta y = y_2 - y_1 = f(x_2) - f(x_1) = f(x_1 + \Delta x) - f(x_1)
  \]
  \textbf{Tasa de cambio media (velocidad media).}
  \[
    T = \frac{\Delta y}{\Delta x} = \frac{f(x_2) - f(x_1)}{x_2 - x_1} = \frac{f(x_1 + \Delta x) - f(x_1)}{\Delta x}
  \]
\end{cuadroteoria}

\begin{cuadropasos}
  \textbf{Cómo calcular una tasa de cambio media:}
  \begin{enumerate}
    \item Identifica $x_1$ y $x_2$ y calcula $\Delta x = x_2 - x_1$.
    \item Evalúa $y_1 = f(x_1)$ y $y_2 = f(x_2)$.
    \item Calcula $\Delta y = y_2 - y_1$.
    \item Divide: $T = \Delta y / \Delta x$.
  \end{enumerate}
\end{cuadropasos}

\begin{cuadroadvertencia}
  \textbf{Ojo con esto:} $\Delta x$ es un \emph{incremento}, no una multiplicación. Puede ser \textbf{negativo} (si $x_2 < x_1$) y siempre se calcula como \emph{valor final menos valor inicial}. Un $\Delta y$ distinto de cero indica cambio; la tasa $T$ tiene unidades de $y$ por unidad de $x$.
\end{cuadroadvertencia}

% ─────────────────────────────────────────────────────────────
\subsection{Tema 2 · Modelos matemáticos en economía}

En economía, muchas magnitudes se relacionan mediante funciones. Las más importantes son el ingreso, el costo y el beneficio.

\begin{cuadroteoria}
  \textbf{Funciones económicas básicas} ($q$ $=$ unidades producidas, $P$ $=$ precio):
  \begin{itemize}
    \item \textbf{Ingreso:} $I(q) = P(q)\cdot q$ \qquad (forma simple: $I(q) = q\,P$)
    \item \textbf{Costo total:} $C(q) = C_v + C_f$ \hfill ($C_v$ $=$ costo variable, $C_f$ $=$ costo fijo)
    \item \textbf{Beneficio:} $B(q) = I(q) - C(q)$
  \end{itemize}
  \textbf{Equilibrio de mercado.} La oferta y la demanda se igualan en el \emph{punto de equilibrio}:
  \[
    S(P) = D(P)
  \]
  \textbf{Interés compuesto.}
  \[
    P = P_0\,(1 + r)^{n}
  \]
  donde $P_0$ es el capital inicial, $r$ la tasa de interés por período y $n$ el número de períodos.
\end{cuadroteoria}

\begin{cuadroextra}
  \textbf{En el equilibrio del productor.} El productor encuentra su punto de equilibrio cuando el ingreso iguala al costo: $I(q) = C(q)$. Es la intersección de ambas curvas y suele dar una ecuación cuadrática con \emph{dos} soluciones; se elige la que tenga sentido económico (la mayor, cuando $I$ crece más rápido que $C$).
\end{cuadroextra}

% ─────────────────────────────────────────────────────────────
\subsection{Tema 3 · La derivada: definición e interpretación}

\begin{cuadroteoria}
  \textbf{De la tasa media a la tasa instantánea.} Si en $T = \dfrac{\Delta y}{\Delta x}$ hacemos que $\Delta x \to 0$, la \textbf{recta secante} que pasa por $P$ y $Q$ tiende a la \textbf{recta tangente} en $P$. Ese límite es la \textbf{derivada}:
  \[
    f'(x) \;=\; \lim_{\Delta x \to 0} \frac{\Delta y}{\Delta x}
    \;=\; \lim_{\Delta x \to 0} \frac{y_2 - y_1}{\Delta x}
    \;=\; \lim_{\Delta x \to 0} \frac{f(x + \Delta x) - f(x)}{\Delta x}
  \]
  \textbf{Condición:} $f$ debe ser \textbf{continua} (y el límite debe existir) para que $f$ sea derivable en $x$.
\end{cuadroteoria}

\begin{cuadroteoria}
  \textbf{Interpretación de la derivada:}
  \begin{itemize}
    \item \textbf{Analítica:} ritmo (razón) de cambio de $f$ respecto a su variable independiente.
    \item \textbf{Geométrica:} $f'(x_0)$ es la \textbf{pendiente de la recta tangente} a la gráfica de $f$ en $x_0$ $\bigl(m_T = f'(x_0)\bigr)$.
    \item \textbf{Nombres equivalentes:} razón de cambio $\cdot$ tasa instantánea $\cdot$ velocidad instantánea.
  \end{itemize}
  \textbf{Notaciones equivalentes de la derivada:}
  \[
    f'(x) \qquad y' \qquad \frac{dy}{dx} \qquad \left(\frac{d}{dx}\right)y
  \]
\end{cuadroteoria}

% ─────────────────────────────────────────────────────────────
\subsection{Tema 4 · Reglas de derivación y técnicas avanzadas}

\begin{cuadroteoria}
  \textbf{Primeras reglas (tabla de derivadas):}
  \begin{center}
  \begin{tabular}{@{}ll@{\qquad}ll@{}}
    \toprule
    \textbf{Función} & \textbf{Derivada} & \textbf{Función} & \textbf{Derivada} \\
    \midrule
    $f = u + v + w$ & $f' = u' + v' + w'$ & $f(x) = ax$ & $f'(x) = a$ \\
    $f = k$ (cte.) & $f' = 0$ & $f(x) = x^{n}$ & $f'(x) = n\,x^{n-1}$ \\
    \bottomrule
  \end{tabular}
  \end{center}
  \textbf{Reglas de producto, cociente y cadena:}
  \[
    f = u\,v \;\Rightarrow\; f' = u'v + uv'
    \qquad
    f = \frac{u}{v} \;\Rightarrow\; f' = \frac{v\,u' - u\,v'}{v^{2}}
    \qquad
    f = u^{n} \;\Rightarrow\; f' = n\,u^{n-1}\,u'
  \]
  \textbf{Exponencial y logaritmo:}
  \[
    f = e^{u} \;\Rightarrow\; f' = e^{u}\,u'
    \qquad
    f = \ln u \;\Rightarrow\; f' = \frac{u'}{u}
  \]
\end{cuadroteoria}

\begin{cuadropasos}
  \textbf{Derivación logarítmica} (para funciones de la forma $y = u(x)^{v(x)}$):
  \begin{enumerate}
    \item Aplica $\ln$ en ambos lados: $\ln y = \ln\!\bigl[u^{v}\bigr]$.
    \item Baja el exponente con $\log a^{n} = n\log a$: $\ln y = v\,\ln u$.
    \item Deriva implícitamente usando $\dfrac{d}{dx}[\ln y] = \dfrac{y'}{y}$ y la regla del producto.
    \item Despeja $y' = y\,[\;\cdots\;]$ y sustituye $y = u^{v}$.
  \end{enumerate}
\end{cuadropasos}

\begin{cuadroteoria}
  \textbf{Derivación implícita.} Si $x$ e $y$ están ligadas por una ecuación $F(x,y)=0$, se deriva término a término respecto a $x$ tratando a $y$ como función de $x$ (cada vez que se deriva $y$ aparece un factor $y'$). Luego se agrupan los términos con $y'$ y se despeja.
  \\[6pt]
  \textbf{Derivadas sucesivas:}
  \begin{center}
  \begin{tabular}{@{}ll@{}}
    \toprule
    \textbf{Orden} & \textbf{Notaciones} \\
    \midrule
    Primera derivada  & $y'$ \quad $f'(x)$ \quad $\dfrac{dy}{dx}$ \\
    Segunda derivada  & $y''$ \quad $f''(x)$ \quad $\dfrac{d^{2}y}{dx^{2}}$ \\
    Tercera derivada  & $y'''$ \quad $f'''(x)$ \quad $\dfrac{d^{3}y}{dx^{3}}$ \\
    $n$-ésima derivada & $\dfrac{d^{n}y}{dx^{n}}$ \\
    \bottomrule
  \end{tabular}
  \end{center}
\end{cuadroteoria}

\begin{cuadroextra}
  \textbf{En economía: las derivadas marginales.}
  \begin{center}
  \begin{tabular}{@{}llll@{}}
    \toprule
    \textbf{Magnitud} & \textbf{Función} & \textbf{Derivada} & \textbf{Nombre} \\
    \midrule
    Ingreso  & $I(q)$ & $I'(q) = \dfrac{dI}{dq}$ & \textbf{Ingreso marginal} \\
    Costo    & $C(q)$ & $C'(q) = \dfrac{dC}{dq}$ & \textbf{Costo marginal} \\
    Beneficio & $B(q)$ & $B'(q) = \dfrac{dB}{dq} = \dfrac{d}{dq}(I - C)$ & \textbf{Beneficio marginal} \\
    \bottomrule
  \end{tabular}
  \end{center}
\end{cuadroextra}

\newpage
% ============================================================
\section{Ejemplos Resueltos}
% ============================================================

\noindent\textbf{Ejemplo resuelto 1 (incrementos y tasa de cambio).} \emph{[Clase 1]}
Hallar los incrementos de $x$ y de $y$ para $x_1 = 2$ y $x_2 = 2{,}03$ en la función $F(x) = \sqrt{3x^{2} + 1}$.

\begin{cuadrosolucion}
  \textbf{Paso 1: incremento de $x$.}
  \[
    \Delta x = x_2 - x_1 = 2{,}03 - 2 = 0{,}03
  \]
  \textbf{Paso 2: valores de la función.}
  \begin{align*}
    y_1 &= F(2) = \sqrt{3(2)^{2}+1} = \sqrt{13} \approx 3{,}61 \\
    y_2 &= F(2{,}03) = \sqrt{3(2{,}03)^{2}+1} = \sqrt{13{,}3627} \approx 3{,}66
  \end{align*}
  \textbf{Paso 3: incremento de $y$.}
  \[
    \Delta y = y_2 - y_1 \approx 3{,}66 - 3{,}61 = 0{,}05
  \]
  \textbf{Paso 4: tasa de cambio media.}
  \[
    T = \frac{\Delta y}{\Delta x} = \frac{0{,}05}{0{,}03} = \frac{5}{3} \approx 1{,}67
  \]
  \textbf{Respuesta:} $\Delta x = 0{,}03$, $\Delta y \approx 0{,}05$ y $T \approx 1{,}67$.
\end{cuadrosolucion}

\noindent\textbf{Ejemplo resuelto 2 (equilibrio de oferta y demanda).} \emph{[Clase 3]}
Las funciones de oferta y demanda de un artículo son $S(P) = 4P + 200$ y $D(P) = -3P + 480$. Hallar el precio de equilibrio y las cantidades correspondientes. Si el precio se incrementa un $2\%$ a partir del equilibrio, ¿cuál es el incremento de la oferta y de la demanda?

\begin{cuadrosolucion}
  \textbf{Paso 1: condición de equilibrio} $S(P) = D(P)$.
  \[
    4P + 200 = -3P + 480
    \;\Longrightarrow\; 7P = 280
    \;\Longrightarrow\; \boxed{P = 40}
  \]
  \textbf{Paso 2: cantidades en el equilibrio.}
  \[
    S(40) = 4(40) + 200 = 360
    \qquad
    D(40) = -3(40) + 480 = 360
  \]
  \[
    \Longrightarrow\quad \text{Punto de equilibrio } (40,\;360)
  \]
  \textbf{Paso 3: incremento del $2\%$ en el precio.}
  \[
    40 \to 100\%, \quad x \leftarrow 2\%
    \;\Longrightarrow\; \Delta P = \frac{40 \cdot 2}{100} = 0{,}8
    \;\Longrightarrow\; P_2 = 40{,}8
  \]
  \textbf{Paso 4: incremento de la oferta.}
  \[
    S(40{,}8) = 4(40{,}8) + 200 = 363{,}2
    \;\Longrightarrow\; \Delta S = 363{,}2 - 360 = 3{,}2
  \]
  \textbf{Paso 5: incremento de la demanda.}
  \[
    D(40{,}8) = -3(40{,}8) + 480 = 357{,}6
    \;\Longrightarrow\; \Delta D = 357{,}6 - 360 = -2{,}4
  \]
  \textbf{Paso 6: tasas de cambio.}
  \[
    T_S = \frac{\Delta S}{\Delta P} = \frac{3{,}2}{0{,}8} = 4
    \qquad
    T_D = \frac{\Delta D}{\Delta P} = \frac{-2{,}4}{0{,}8} = -3
  \]
  \textbf{Respuesta:} $P=40$, $Q=360$; $\Delta S = 3{,}2$, $\Delta D = -2{,}4$; $T_S = 4$ y $T_D = -3$.
\end{cuadrosolucion}

\noindent\textbf{Ejemplo resuelto 3 (equilibrio del productor).} \emph{[Clase 4]}
El chocolate Savoy vende cajas de bombones a $2000$ bs c/u. Si $q$ es el número de cajas producidas en miles semanalmente, los costos son $C(q) = 100q^{2} + 1300q + 1000$. Determine el punto de equilibrio del productor. Si $q$ se incrementa un $5\%$ a partir del equilibrio, determine el incremento del costo y del ingreso.

\begin{cuadrosolucion}
  \textbf{Paso 1: función ingreso.}
  \[
    I(q) = q\,P = 2000q
  \]
  \textbf{Paso 2: equilibrio del productor} $I(q) = C(q)$.
  \[
    2000q = 100q^{2} + 1300q + 1000
    \;\Longrightarrow\; 100q^{2} - 700q + 1000 = 0
  \]
  Dividiendo entre $100$: $\;q^{2} - 7q + 10 = 0$, con $a=1$, $b=-7$, $c=10$:
  \[
    q = \frac{7 \pm \sqrt{49 - 40}}{2} = \frac{7 \pm 3}{2}
    \;\Longrightarrow\; q_1 = 5, \quad q_2 = 2
  \]
  \textbf{Paso 3: verificación.}
  \[
    I(5) = 10000 = C(5)
    \qquad
    I(2) = 4000 = C(2)
  \]
  Se toma $q_1 = 5$ (la cantidad que maximiza el beneficio).
  \textbf{Paso 4: incremento del $5\%$ en $q$.}
  \[
    5 \to 100\%, \quad x \leftarrow 5\%
    \;\Longrightarrow\; \Delta q = \frac{5 \cdot 5}{100} = 0{,}25
    \;\Longrightarrow\; q_2 = 5{,}25
  \]
  \textbf{Paso 5: incremento del ingreso.}
  \[
    I_1 = 2000(5) = 10\,000
    \qquad
    I_2 = 2000(5{,}25) = 10\,500
    \;\Longrightarrow\; \Delta I = 10\,500 - 10\,000 = 500
  \]
  \textbf{Paso 6: incremento del costo.}
  \[
    C_1 = 100(5)^{2} + 1300(5) + 1000 = 10\,000
  \]
  \[
    C_2 = 100(5{,}25)^{2} + 1300(5{,}25) + 1000
        = 2756{,}25 + 6825 + 1000 = 10\,581{,}25
  \]
  \[
    \Delta C = 10\,581{,}25 - 10\,000 = 581{,}25
  \]
  \textbf{Respuesta:} equilibrio en $q=5$; $\Delta I = 500$ y $\Delta C = 581{,}25$.
  \\[2pt]
  \textit{Fe de erratas:} el cuaderno registra $C_2 = 10\,256{,}25$ y $\Delta C = 256{,}25$; el error proviene de usar $1300 \cdot 5$ en lugar de $1300 \cdot 5{,}25$.
\end{cuadrosolucion}

\noindent\textbf{Ejemplo resuelto 4 (interés compuesto y logaritmos).} \emph{[Clase 5]}
Si $\$100$ se invierten a un interés compuesto de $6\%$ anual, ¿cuánto tardará la inversión en llegar a $\$150$?

\begin{cuadrosolucion}
  \textbf{Datos:} $P_0 = 100$, $P_F = 150$, $r = 6\% = 0{,}06$, $n = ?$
  \[
    P_F = P_0(1+r)^{n}
    \;\Longrightarrow\; 100(1{,}06)^{n} = 150
    \;\Longrightarrow\; (1{,}06)^{n} = 1{,}5
  \]
  \textbf{Aplicamos logaritmos:}
  \[
    \log(1{,}06)^{n} = \log(1{,}5)
    \;\Longrightarrow\; n\,\log(1{,}06) = \log(1{,}5)
    \;\Longrightarrow\; n = \frac{\log 1{,}5}{\log 1{,}06} \approx \frac{0{,}1761}{0{,}0253} \approx 6{,}96
  \]
  \textbf{Respuesta:} $n \approx 7$ años.
\end{cuadrosolucion}

\noindent\textbf{Ejemplo resuelto 5 (derivada por definición, $\Delta x$).} \emph{[Clase 7.1]}
Hallar la derivada de $f(x) = 3x^{2} - 2x + 1$ por definición.

\begin{cuadrosolucion}
  \textbf{Paso 1:} planteamos el límite del cociente incremental con $\Delta x$.
  \[
    f'(x) = \lim_{\Delta x \to 0} \frac{f(x + \Delta x) - f(x)}{\Delta x}
  \]
  \textbf{Paso 2:} sustituimos y desarrollamos el numerador.
  \[
    = \lim_{\Delta x \to 0} \frac{\bigl[3(x+\Delta x)^{2} - 2(x+\Delta x) + 1\bigr] - \bigl[3x^{2} - 2x + 1\bigr]}{\Delta x}
  \]
  \[
    = \lim_{\Delta x \to 0} \frac{3x^{2} + 6x\,\Delta x + 3(\Delta x)^{2} - 2x - 2\Delta x + 1 - 3x^{2} + 2x - 1}{\Delta x}
  \]
  \textbf{Paso 3:} cancelamos $3x^{2}$, $2x$ y $1$.
  \[
    = \lim_{\Delta x \to 0} \frac{6x\,\Delta x + 3(\Delta x)^{2} - 2\Delta x}{\Delta x}
    = \lim_{\Delta x \to 0} \frac{\Delta x\,\bigl(6x + 3\Delta x - 2\bigr)}{\Delta x}
  \]
  \textbf{Paso 4:} simplificamos $\Delta x$ y tomamos el límite.
  \[
    f'(x) = \lim_{\Delta x \to 0} \bigl(6x + 3\Delta x - 2\bigr) = 6x + 3(0) - 2 = 6x - 2
  \]
  \textbf{Respuesta:} $f'(x) = 6x - 2$ \quad \textit{(coincide con la tabla: $2\cdot 3x - 2$).}
\end{cuadrosolucion}

\noindent\textbf{Ejemplo resuelto 6 (derivada por definición de un cociente, $\Delta x$).} \emph{[Clase 7.2]}
Hallar la derivada de $g(x) = \dfrac{1}{x+2}$ por definición.

\begin{cuadrosolucion}
  \textbf{Paso 1:} planteamos el límite.
  \[
    g'(x) = \lim_{\Delta x \to 0} \frac{g(x+\Delta x) - g(x)}{\Delta x}
    = \lim_{\Delta x \to 0} \frac{\dfrac{1}{x + \Delta x + 2} - \dfrac{1}{x+2}}{\Delta x}
  \]
  \textbf{Paso 2:} común denominador $(x+\Delta x+2)(x+2)$.
  \[
    = \lim_{\Delta x \to 0} \frac{(x+2) - (x + \Delta x + 2)}{(x+\Delta x+2)(x+2)\,\Delta x}
    = \lim_{\Delta x \to 0} \frac{-\Delta x}{(x+\Delta x+2)(x+2)\,\Delta x}
  \]
  \textbf{Paso 3:} simplificamos $\Delta x$ y tomamos el límite.
  \[
    = \lim_{\Delta x \to 0} \frac{-1}{(x+\Delta x+2)(x+2)}
    = \frac{-1}{(x+0+2)(x+2)}
  \]
  \textbf{Respuesta:} $g'(x) = -\dfrac{1}{(x+2)^{2}}$.
\end{cuadrosolucion}

\newpage
\noindent\textbf{Ejemplo resuelto 7 (regla del cociente).} \emph{[Clase 10.2]}
Derivar $h(x) = \dfrac{4x^{2}+1}{3x+2}$.

\begin{cuadrosolucion}
  \textbf{Paso 1:} identificamos $u = 4x^{2}+1$ y $v = 3x+2$.
  \[
    u' = 8x
    \qquad
    v' = 3
  \]
  \textbf{Paso 2:} aplicamos $f' = \dfrac{v\,u' - u\,v'}{v^{2}}$.
  \[
    h'(x) = \frac{(3x+2)(8x) - (4x^{2}+1)(3)}{(3x+2)^{2}}
  \]
  \textbf{Paso 3:} desarrollamos y simplificamos el numerador.
  \[
    h'(x) = \frac{24x^{2} + 16x - 12x^{2} - 3}{(3x+2)^{2}}
  \]
  \textbf{Respuesta:} $h'(x) = \dfrac{12x^{2} + 16x - 3}{(3x+2)^{2}}$.
\end{cuadrosolucion}

\noindent\textbf{Ejemplo resuelto 8 (regla del producto).} \emph{[Clase 11.1]}
Derivar $F(t) = (2t^{3}+t^{2})(5t+1)$.

\begin{cuadrosolucion}
  \textbf{Paso 1:} identificamos $u = 2t^{3}+t^{2}$ y $v = 5t+1$.
  \[
    u' = 6t^{2}+2t
    \qquad
    v' = 5
  \]
  \textbf{Paso 2:} aplicamos $f' = u'v + uv'$.
  \[
    F'(t) = (6t^{2}+2t)(5t+1) + (2t^{3}+t^{2})(5)
  \]
  \textbf{Paso 3:} desarrollamos y agrupamos términos semejantes.
  \[
    F'(t) = 30t^{3} + 6t^{2} + 10t^{2} + 2t + 10t^{3} + 5t^{2}
  \]
  \textbf{Respuesta:} $F'(t) = 40t^{3} + 21t^{2} + 2t$.
\end{cuadrosolucion}

\noindent\textbf{Ejemplo resuelto 9 (regla de la cadena con raíz).} \emph{[Clase 11.2]}
Derivar $F(x) = \sqrt{3x^{2}+5x-1}$.

\begin{cuadrosolucion}
  \textbf{Paso 1:} reescribimos con exponente fraccionario.
  \[
    F(x) = (3x^{2}+5x-1)^{1/2}
  \]
  \textbf{Paso 2:} aplicamos $f' = \tfrac{1}{2}u^{-1/2}\,u'$ con $u = 3x^{2}+5x-1$ y $u' = 6x+5$.
  \[
    F'(x) = \frac{1}{2}(3x^{2}+5x-1)^{-1/2}\,(6x+5)
  \]
  \textbf{Respuesta:} $F'(x) = \dfrac{6x+5}{2\sqrt{3x^{2}+5x-1}}$.
\end{cuadrosolucion}

\noindent\textbf{Ejemplo resuelto 10 (exponencial y logaritmo).} \emph{[Clase 12.2–12.3]}

\begin{cuadrosolucion}
  \textbf{(a) Exponencial:} $y = e^{3x^{2}+1}$. Con $u = 3x^{2}+1$ se tiene $u' = 6x$:
  \[
    y' = \frac{dy}{dx} = e^{3x^{2}+1}\,(6x) = 6x\,e^{3x^{2}+1}
  \]
  \textbf{(b) Logaritmo:} $P(x) = \ln(2x^{2}+x)$. Con $u = 2x^{2}+x$ se tiene $u' = 4x+1$:
  \[
    P'(x) = \frac{4x+1}{2x^{2}+x}
  \]
\end{cuadrosolucion}

\noindent\textbf{Ejemplo resuelto 11 (derivación logarítmica).} \emph{[Clase 15.1]}
Derivar $y = (5x + e^{x})^{2x+1}$.

\begin{cuadrosolucion}
  \textbf{Paso 1:} aplicamos $\ln$ en ambos lados.
  \[
    \ln y = \ln (5x + e^{x})^{2x+1}
  \]
  \textbf{Paso 2:} bajamos el exponente con $\log a^{n} = n\log a$.
  \[
    \ln y = (2x+1)\,\ln(5x+e^{x})
  \]
  \textbf{Paso 3:} derivamos con la regla del producto y $\dfrac{d}{dx}[\ln y] = \dfrac{y'}{y}$.
  \[
    \frac{y'}{y} = (2x+1)\cdot\frac{5 + e^{x}}{5x + e^{x}} + \ln(5x+e^{x})\cdot 2
  \]
  \textbf{Paso 4:} despejamos $y'$ y sustituimos $y$.
  \[
    y' = (5x+e^{x})^{2x+1}\left[\frac{(2x+1)(5+e^{x})}{5x+e^{x}} + 2\ln(5x+e^{x})\right]
  \]
  \textbf{Respuesta:} $y' = (5x+e^{x})^{2x+1}\left[\dfrac{(2x+1)(5+e^{x})}{5x+e^{x}} + 2\ln(5x+e^{x})\right]$.
\end{cuadrosolucion}

\newpage
\noindent\textbf{Ejemplo resuelto 12 (derivación implícita).} \emph{[Clase 16.1]}
Hallar $y'$ si $a\,x^{3} + 4x^{2}y - x\,y^{3} = 0$.

\begin{cuadrosolucion}
  \textbf{Paso 1:} derivamos término a término respecto a $x$ (cada derivada de $y$ lleva factor $y'$).
  \[
    3a\,x^{2} + \bigl(8xy + 4x^{2}y'\bigr) - \bigl(y^{3} + x\cdot 3y^{2}y'\bigr) = 0
  \]
  \textbf{Paso 2:} agrupamos los términos con $y'$.
  \[
    3a\,x^{2} + 8xy + 4x^{2}y' - y^{3} - 3xy^{2}y' = 0
    \;\Longrightarrow\;
    y'\bigl(4x^{2} - 3xy^{2}\bigr) = y^{3} - 3a\,x^{2} - 8xy
  \]
  \textbf{Respuesta:} $y' = \dfrac{y^{3} - 3a\,x^{2} - 8xy}{4x^{2} - 3xy^{2}}$.
\end{cuadrosolucion}

\noindent\textbf{Ejemplo resuelto 13 (derivadas sucesivas y demostración).} \emph{[Clase 17.3]}
Dada $y = A\,e^{x} + B\,e^{x/2}$, demostrar que $2y'' - 3y' + y = 0$.

\begin{cuadrosolucion}
  \textbf{Paso 1:} calculamos las derivadas.
  \[
    y = A\,e^{x} + B\,e^{x/2}
    \qquad
    y' = A\,e^{x} + \tfrac{1}{2}B\,e^{x/2}
    \qquad
    y'' = A\,e^{x} + \tfrac{1}{4}B\,e^{x/2}
  \]
  \textbf{Paso 2:} sustituimos en $2y'' - 3y' + y$.
  \[
    2\left(Ae^{x} + \tfrac{1}{4}Be^{x/2}\right) - 3\left(Ae^{x} + \tfrac{1}{2}Be^{x/2}\right) + Ae^{x} + Be^{x/2}
  \]
  \[
    = 2Ae^{x} + \tfrac{1}{2}Be^{x/2} - 3Ae^{x} - \tfrac{3}{2}Be^{x/2} + Ae^{x} + Be^{x/2}
  \]
  \textbf{Paso 3:} agrupamos por $e^{x}$ y $e^{x/2}$.
  \[
    = Ae^{x}\,(2 - 3 + 1) + Be^{x/2}\left(\tfrac{1}{2} - \tfrac{3}{2} + 1\right)
    = Ae^{x}\cdot 0 + Be^{x/2}\cdot 0 = 0
  \]
  \textbf{Respuesta:} $2y'' - 3y' + y = 0$. \hfill $\blacksquare$
\end{cuadrosolucion}

\newpage
% ============================================================
\section{Ejercicios Propuestos}
% ============================================================

\subsection{Ejercicios nuevos}

\noindent En los ejercicios 1--6 $(\bigstar\bigstar)$, 7--13 $(\bigstar\bigstar\bigstar)$ y 14--18 $(\bigstar\bigstar\bigstar\bigstar)$ la dificultad aumenta progresivamente.

\begin{enumerate}[leftmargin=*]

  % ── ★★ (1─6) ─────────────────────────────────────────────
  \item \medio{} Para $f(x) = x^{2} + 3x$ con $x_1 = 1$ y $x_2 = 1{,}04$, calcule $\Delta x$, $\Delta y$ y la tasa de cambio media $T$.

  \item \medio{} Si $\Delta x = -0{,}08$ y $x_1 = 12$, halle $x_2$ (despeje a partir de $\Delta x = x_2 - x_1$).

  \item \medio{} Se invierten $\$500$ a interés compuesto del $4\%$ anual. ¿Cuál es el capital al cabo de $5$ años?

  \item \medio{} Un artículo tiene funciones de oferta y demanda $S(P) = 5P + 120$ y $D(P) = -4P + 390$. Halle el precio y la cantidad de equilibrio.

  \item \medio{} Calcule $f'(x)$ por definición (con $\Delta x$) para $f(x) = x^{2} - 5x$.

  \item \medio{} Derive por la regla del cociente: $f(x) = \dfrac{2x^{2} - 3}{x + 1}$.

  % ── ★★★ (7─13) ───────────────────────────────────────────
  \item \dificil{} Derive por la regla del producto: $f(x) = (x^{2}+1)(3x-2)$.

  \item \dificil{} Derive por la regla de la cadena: $f(x) = (4x^{2}-1)^{5}$.

  \item \dificil{} Derive: $f(x) = \sqrt{5x^{2} + 2x}$.

  \item \dificil{} Derive: $y = e^{\,x^{2} - 3x}$.

  \item \dificil{} Derive: $f(x) = \ln(3x^{2} - x)$.

  \item \dificil{} Derive: $h(x) = \ln\!\left(\dfrac{x^{2}+1}{x-2}\right)$.

  \item \dificil{} Halle $y'$ por derivación implícita: $x^{2} + 4xy - y^{2} = 0$.

  % ── ★★★★ (14─18) ─────────────────────────────────────────
  \item \muyDificil{} Halle $y'$ por derivación logarítmica: $y = x^{x}$.

  \item \muyDificil{} Halle $y'$ por derivación logarítmica: $y = (2x+1)^{3x}$.

  \item \muyDificil{} Halle $y'$ por derivación implícita: $3x^{2} + 2y^{3} - \ln y = 0$.

  \item \muyDificil{} Dada $f(x) = x^{4} - 3x^{3} + 2x$, calcule $f''(x)$ y $f'''(x)$.

  \item \muyDificil{} Dada $y = A\,e^{3x} + B\,e^{x}$, demuestre que $y'' - 4y' + 3y = 0$.

\end{enumerate}

\begin{enumerate}[leftmargin=*, start=19]

  \item \dificil{} \textbf{(Clase)} Derive $g(x) = \sqrt[3]{x^{2}} - 4\sqrt{x^{3}}$.

  \item \dificil{} \textbf{(Clase)} Derive $Z(x) = \dfrac{1}{\sqrt{5x+3}}$.

  \item \muyDificil{} \textbf{(Clase)} Derive $h(x) = \ln\sqrt{\dfrac{x^{2}+3}{2x-1}}$.

  \item \dificil{} \textbf{(Clase)} Derive $f(x) = \ln\!\left[\dfrac{x^{2}-2x+1}{3x-1}\right]$.

  \item \dificil{} \textbf{(Clase)} Derive $g(x) = \sqrt{3x^{2} - x + 1}$.

  \item \dificil{} \textbf{(Clase)} Derive $h(x) = 2x^{p} + x^{p-1}$.

  \item \dificil{} \textbf{(Clase)} Derive $M(x) = e^{-x^{2}} + \ln(2x^{n} - 3)$.

  \item \muyDificil{} \textbf{(Clase)} Derive por derivación logarítmica: $J(x) = \left[\dfrac{3x^{2}-1}{2x+1}\right]^{x+2}$.

  \item \muyDificil{} \textbf{(Clase)} Derive $f(p) = \dfrac{e^{-p^{2}} + \ln(p^{2})}{p^{2}+1}$.

  \item \muyDificil{} \textbf{(Clase)} Derive $Q(s) = \sqrt[5]{25^{s} - e^{3s}}$.

  \item \muyDificil{} \textbf{(Clase)} Derive $h(z) = \ln\!\left(\sqrt{z^{2} + e^{-z}}\right)^{3}$.

  \item \muyDificil{} \textbf{(Clase)} Derive $W(x) = \dfrac{3x^{2}\,e^{2x}}{x-1}$.

  \item \muyDificil{} \textbf{(Clase)} Halle $y'$ por derivación implícita: $3x^{2} + 6y^{3} - \ln y = 0$.

  \item \muyDificil{} \textbf{(Clase)} Halle $y'$ por derivación implícita: $y^{2} - 3e^{-x} + \ln y = \ln(3x+1)$.

  \item \muyDificil{} \textbf{(Clase)} Halle $y'$ por derivación implícita: $\sqrt{y^{2}-1} + 3xy + \sqrt{3x-2} = 0$.

\end{enumerate}

\newpage
% ============================================================
\ifmwclaves
  \section{Clave de Respuestas}
  % ============================================================

  \begin{enumerate}[leftmargin=*, label=\textbf{\arabic*.}]

    % ── Nuevos 1─18 ─────────────────────────────────────────
    \item $\Delta x = 0{,}04$; \quad $\Delta y = 0{,}2016$; \quad $T = 5{,}04$

    \item $x_2 = x_1 + \Delta x = 12 - 0{,}08 = 11{,}92$

    \item $P = 500(1{,}04)^{5} \approx \$608{,}33$

    \item $P = 30$; \quad $Q = S(30) = D(30) = 270$

    \item $f'(x) = 2x - 5$

    \item $f'(x) = \dfrac{2x^{2} + 4x + 3}{(x+1)^{2}}$

    \item $f'(x) = 9x^{2} - 4x + 3$

    \item $f'(x) = 40x\,(4x^{2}-1)^{4}$

    \item $f'(x) = \dfrac{5x+1}{\sqrt{5x^{2}+2x}}$

    \item $y' = (2x-3)\,e^{\,x^{2}-3x}$

    \item $f'(x) = \dfrac{6x-1}{3x^{2}-x}$

    \item $h'(x) = \dfrac{x^{2} - 4x - 1}{(x^{2}+1)(x-2)}$

    \item $y' = -\dfrac{x + 2y}{2x - y}$

    \item $y' = x^{x}\,(1 + \ln x)$

    \item $y' = (2x+1)^{3x}\left[3\ln(2x+1) + \dfrac{6x}{2x+1}\right]$

    \item $y' = -\dfrac{6xy}{6y^{3} - 1}$

    \item $f''(x) = 12x^{2} - 18x$; \quad $f'''(x) = 24x - 18$

    \item Demostración: $y' = 3Ae^{3x} + Be^{x}$, \; $y'' = 9Ae^{3x} + Be^{x}$; \; sustituyendo, $y'' - 4y' + 3y = A(9-12+3)e^{3x} + B(1-4+3)e^{x} = 0$.

    % ── De clase 19─33 ──────────────────────────────────────
    \item $g'(x) = \dfrac{2}{3\sqrt[3]{x}} - 6\sqrt{x}$

    \item $Z'(x) = -\dfrac{5}{2\,(5x+3)^{3/2}}$

    \item $h'(x) = \dfrac{1}{2}\cdot\dfrac{2x^{2} - 2x - 6}{(x^{2}+3)(2x-1)} = \dfrac{x^{2} - x - 3}{(x^{2}+3)(2x-1)}$

    \item $f'(x) = \dfrac{2x-2}{x^{2}-2x+1} - \dfrac{3}{3x-1}$

    \item $g'(x) = \dfrac{6x-1}{2\sqrt{3x^{2}-x+1}}$

    \item $h'(x) = 2p\,x^{p-1} + (p-1)\,x^{p-2}$

    \item $M'(x) = -2x\,e^{-x^{2}} + \dfrac{2n\,x^{n-1}}{2x^{n}-3}$

    \item $J'(x) = \left[\dfrac{3x^{2}-1}{2x+1}\right]^{x+2}\left[(x+2)\left(\dfrac{6x}{3x^{2}-1} - \dfrac{2}{2x+1}\right) + \ln(3x^{2}-1) - \ln(2x+1)\right]$

    \item $f'(p) = \dfrac{-2p\,(p^{2}+2)\,e^{-p^{2}} + \dfrac{2(p^{2}+1)}{p} - 4p\ln p}{(p^{2}+1)^{2}}$

    \item $Q'(s) = \dfrac{1}{5}\left(25^{s} - e^{3s}\right)^{-4/5}\left(25^{s}\ln 25 - 3e^{3s}\right)$

    \item $h'(z) = \dfrac{3}{2}\cdot\dfrac{2z - e^{-z}}{z^{2} + e^{-z}}$

    \item $W'(x) = \dfrac{3x\,e^{2x}\,(2x^{2} - x - 2)}{(x-1)^{2}}$

    \item $y' = -\dfrac{6xy}{18y^{3} - 1}$

    \item $y' = \dfrac{\dfrac{3}{3x+1} - 3e^{-x}}{2y + \dfrac{1}{y}}$

    \item $y' = \dfrac{-3y - \dfrac{3}{2\sqrt{3x-2}}}{y\,(y^{2}-1)^{-1/2} + 3x}$

  \end{enumerate}
\fi

\end{document}
